\documentclass[12pt]{article}

\usepackage[margin=1in]{geometry}
\usepackage{amsmath, amssymb, amsthm}
\usepackage[T1]{fontenc}
\usepackage{lmodern}
\usepackage{microtype}
\usepackage{setspace}
\usepackage{booktabs}
\usepackage{array}
\usepackage{caption}
\usepackage[colorlinks=true, allcolors=blue]{hyperref}

\setlength{\parindent}{0pt}
\setlength{\parskip}{0.6em}
\onehalfspacing

\captionsetup{
    font=small,
    labelfont=bf,
    justification=centering
}

\newtheorem{theorem}{Theorem}[section]
\newtheorem{lemma}[theorem]{Lemma}
\newtheorem{proposition}[theorem]{Proposition}
\newtheorem{corollary}[theorem]{Corollary}

\theoremstyle{definition}
\newtheorem{definition}[theorem]{Definition}

\theoremstyle{remark}
\newtheorem{remark}[theorem]{Remark}

\title{\textbf{Scott's Factorial Series}}
\author{Jaxson Scottlevoy Freeman\\
\small Independent student researcher\\
\small Corresponding author: \texttt{jaxsonstormy@gmail.com}}
\date{October 2026}

\begin{document}

\maketitle

\begin{abstract}
This paper studies the function defined by the series
\[
S(x)=
\sum_{n=0}^{\infty}
\frac{(-1)^n}{e^n(n!)^2}
\left[
\sin\left(\frac{x}{(n+1)!}\right)
\right]^{2n+1}.
\]
The series combines alternating signs, factorial coefficients,
and increasing odd powers of the sine function. We prove that
it converges absolutely and uniformly on the real line. By
extending the analysis to complex inputs, we show that the
resulting function is entire and therefore has derivatives
of every order. We also derive its Taylor expansion near the
origin, establish global bounds, and obtain an error estimate
for finite approximations. Finally, we examine the periodicity
of the partial sums and discuss possible generalizations.
\end{abstract}

\noindent\textbf{Keywords:} factorial series; uniform convergence; entire functions; Taylor expansion; approximation error.

\medskip
\noindent\textbf{2020 Mathematics Subject Classification:} Primary 30D10; Secondary 40A30, 26A24.

\section{Introduction}

Infinite series provide a useful way to define functions and
study their properties. The behavior of a series often depends
on the relationship between its coefficients and its individual
terms. Factorial coefficients are particularly useful because
they can force rapid convergence even when the remaining
factors vary with the input.

Consider the function
\[
S(x)=
\sum_{n=0}^{\infty}
\frac{(-1)^n}{e^n(n!)^2}
\left[
\sin\left(\frac{x}{(n+1)!}\right)
\right]^{2n+1}.
\]

Three features distinguish this series. Its coefficients
decrease factorially, the argument of the sine contains a
factorial, and the exponent of the sine increases with the
summation index. The signs also alternate.

These features raise several questions. Does the series
converge for every real input? Is its sum continuous and
differentiable? Can the function be represented by a Taylor
series, and how accurately can it be approximated by a finite
sum?

The first part of this paper establishes convergence and
continuity. The subsequent sections examine stronger
regularity properties, behavior near the origin, global
bounds, and numerical approximation. The arguments use
standard results on uniform convergence and holomorphic
functions; see, for example, Rudin~\cite{rudin} and
Conway~\cite{conway}.

\section{Definition and Basic Properties}

\subsection{The series and its partial sums}

\begin{definition}[Scottlevoy's Factorial Series]
For \(x\in\mathbb{R}\), define
\[
S(x)=
\sum_{n=0}^{\infty}
\frac{(-1)^n}{e^n(n!)^2}
\left[
\sin\left(\frac{x}{(n+1)!}\right)
\right]^{2n+1}.
\]
\end{definition}

Here, \(e\) is Euler's number and \(n!\) denotes the factorial
of the nonnegative integer \(n\), with \(0!=1\).

Write the \(n\)-th term as
\[
a_n(x)=
\frac{(-1)^n}{e^n(n!)^2}
\left[
\sin\left(\frac{x}{(n+1)!}\right)
\right]^{2n+1}.
\]

The \(N\)-th partial sum is
\begin{equation}
S_N(x)=\sum_{n=0}^{N}a_n(x).
\label{eq:partial-sum}
\end{equation}

Each partial sum is a finite combination of trigonometric
functions. The infinite series is obtained by taking the
limit of these partial sums, provided that the limit exists.

\subsection{The first terms}

The first five terms are displayed in
Table~\ref{tab:first-terms}.

\begin{table}[htbp]
\centering
\caption{The first five terms of the series.}
\label{tab:first-terms}
\renewcommand{\arraystretch}{1.5}
\setlength{\tabcolsep}{7pt}

\begin{tabular}{c c c c}
\toprule
\(n\) & Sign & Coefficient & Sine factor \\
\midrule
\(0\) & \(+\) & \(1\) &
\(\sin(x)\) \\

\(1\) & \(-\) & \(\dfrac{1}{e}\) &
\(\sin^3\left(\dfrac{x}{2}\right)\) \\

\(2\) & \(+\) & \(\dfrac{1}{4e^2}\) &
\(\sin^5\left(\dfrac{x}{6}\right)\) \\

\(3\) & \(-\) & \(\dfrac{1}{36e^3}\) &
\(\sin^7\left(\dfrac{x}{24}\right)\) \\

\(4\) & \(+\) & \(\dfrac{1}{576e^4}\) &
\(\sin^9\left(\dfrac{x}{120}\right)\) \\
\bottomrule
\end{tabular}
\end{table}

The factorials enter the series in two different ways. The
coefficient decreases as \(n\) increases, while the argument
of the sine becomes smaller for any fixed \(x\). The increasing
exponent further changes the contribution of each term.

These effects suggest that the later terms should become
small. A convergence proof is needed, however, to establish
that their infinite sum is well-defined.

\section{Convergence and Continuity}

\subsection{Absolute and uniform convergence}

\begin{theorem}
The series defining \(S(x)\) converges absolutely and uniformly
on \(\mathbb{R}\).
\end{theorem}

\begin{proof}
For every real \(x\),
\[
\left|\sin\left(\frac{x}{(n+1)!}\right)\right|\leq 1.
\]
It follows that
\[
|a_n(x)|\leq \frac{1}{e^n(n!)^2}.
\]

Consider the numerical series
\[
\sum_{n=0}^{\infty}\frac{1}{e^n(n!)^2}.
\]
The ratio of consecutive terms is
\[
\frac{1/[e^{n+1}((n+1)!)^2]}
     {1/[e^n(n!)^2]}
=\frac{1}{e(n+1)^2}.
\]
This ratio tends to zero as \(n\to\infty\), so the numerical
series converges by the ratio test.

The Weierstrass \(M\)-test therefore proves that the original
series converges uniformly on \(\mathbb{R}\)~\cite{rudin}. Since the same
bound applies to the absolute values of its terms, the
convergence is also absolute at every real \(x\).
\end{proof}

\begin{corollary}
The function \(S\) is continuous on \(\mathbb{R}\).
\end{corollary}

\begin{proof}
Each \(a_n(x)\) is continuous, so every partial sum \(S_N(x)\)
is continuous. The uniform convergence of \(S_N\) to \(S\)
then implies that \(S\) is continuous on \(\mathbb{R}\).
\end{proof}

\subsection{A global bound}

The same estimate used in the convergence proof also gives
a bound on the value of the function.

\begin{proposition}
For every \(x\in\mathbb{R}\),
\[
|S(x)|\leq
\sum_{n=0}^{\infty}\frac{1}{e^n(n!)^2}
\leq e^{1/e}.
\]
\end{proposition}

\begin{proof}
By the triangle inequality and the bound on each term,
\[
|S(x)|\leq
\sum_{n=0}^{\infty}|a_n(x)|
\leq
\sum_{n=0}^{\infty}\frac{1}{e^n(n!)^2}.
\]
Since \(n!\geq 1\), we have
\[
\frac{1}{(n!)^2}\leq\frac{1}{n!}.
\]
Consequently,
\[
\sum_{n=0}^{\infty}\frac{1}{e^n(n!)^2}
\leq
\sum_{n=0}^{\infty}\frac{(1/e)^n}{n!}
=e^{1/e}.
\]
This proves the result.
\end{proof}

Thus, although the sine factors depend on \(x\), the function
remains bounded on the entire real line.

\section{Differentiability and Analyticity}

\subsection{The first derivative}

Differentiating an individual term gives
\begin{align}
a_n'(x)
={}&
\frac{(-1)^n(2n+1)}
{e^n(n!)^2(n+1)!}
\nonumber\\
&\quad{}\times
\sin^{2n}\left(\frac{x}{(n+1)!}\right)
\cos\left(\frac{x}{(n+1)!}\right).
\label{eq:first-derivative}
\end{align}

Because the sine and cosine are bounded in absolute value
by \(1\) on the real line,
\[
|a_n'(x)|
\leq
\frac{2n+1}{e^n(n!)^2(n+1)!}.
\]

The numerical series of these bounds converges by the ratio
test. Therefore, the derivative series converges uniformly
on \(\mathbb{R}\). The theorem on differentiating uniformly
convergent series implies that \(S\) is continuously
differentiable and that
\[
S'(x)=\sum_{n=0}^{\infty}a_n'(x).
\]

This establishes differentiability everywhere, rather than
leaving it as an open question.

\subsection{Higher derivatives}

The regularity of \(S\) can be strengthened further by
considering complex inputs.

\begin{theorem}
The series defining \(S\) converges locally uniformly on
\(\mathbb{C}\). Its sum therefore extends to an entire
function, and all derivatives of \(S\) exist and are
continuous on \(\mathbb{R}\).
\end{theorem}

\begin{proof}
Extend each term to a complex variable \(z\):
\[
a_n(z)=
\frac{(-1)^n}{e^n(n!)^2}
\sin^{2n+1}\left(\frac{z}{(n+1)!}\right).
\]
Each \(a_n\) is entire.

For every complex number \(w\),
\[
|\sin w|\leq e^{|w|}.
\]
Fix \(R>0\) and suppose that \(|z|\leq R\). Then
\[
|a_n(z)|
\leq
\frac{1}{e^n(n!)^2}
\exp\left(\frac{(2n+1)R}{(n+1)!}\right).
\]

The sequence
\[
\frac{2n+1}{(n+1)!}
\]
is bounded above by \(3/2\) for all nonnegative integers
\(n\). Hence
\[
|a_n(z)|
\leq
\frac{e^{3R/2}}{e^n(n!)^2}.
\]

The right-hand side is summable and independent of \(z\)
within the disk \(|z|\leq R\). The Weierstrass \(M\)-test
therefore gives uniform convergence on every closed disk.

A locally uniform limit of entire functions is entire~\cite{conway}.
Consequently, the complex extension of \(S\) is entire,
which proves the result.
\end{proof}

In particular, \(S\) is real analytic: near every real
point, it is represented by a convergent Taylor series.
The Taylor series about the origin in fact converges for
every complex input.

\section{Behavior Near the Origin}

\subsection{The value and derivative at zero}

Since every term contains a positive power of a sine factor,
\[
S(0)=0.
\]

The first term of the series is \(\sin x\). Every remaining
term vanishes to order at least three at the origin.
It follows that
\[
S'(0)=1.
\]

The function therefore behaves like the identity function
near zero, in the sense that
\[
S(x)=x+O(x^3)
\qquad\text{as }x\to 0.
\]

In particular,
\[
\lim_{x\to 0}\frac{S(x)}{x}=1.
\]
This also shows that \(S(x)\) is positive for sufficiently
small positive \(x\) and negative for sufficiently small
negative \(x\).

\subsection{Taylor expansion}

The Taylor expansion of the sine function is
\[
\sin t=t-\frac{t^3}{6}+\frac{t^5}{120}+O(t^7).
\]

For the second term of \(S\), we also use
\[
\sin^3 t=t^3-\frac{t^5}{2}+O(t^7).
\]

The terms with indices \(n=0,1,2\) contribute to the
Taylor expansion through degree five. Every term with
\(n\geq 3\) vanishes to order at least seven at zero.
Combining these contributions gives
\begin{align}
S(x)
={}&x-
\left(\frac{1}{6}+\frac{1}{8e}\right)x^3
\nonumber\\
&+
\left(
\frac{1}{120}
+\frac{1}{64e}
+\frac{1}{31104e^2}
\right)x^5
+O(x^7).
\label{eq:taylor}
\end{align}

The coefficient of \(x^3\), for example, receives
contributions from both \(\sin x\) and
\(-e^{-1}\sin^3(x/2)\). The coefficient of \(x^5\)
receives contributions from the first three terms.

More generally, the term \(a_n(x)\) vanishes to order
\(2n+1\) at zero. Therefore, the coefficient of \(x^{2k+1}\)
in the Taylor expansion receives contributions only from
indices \(n\leq k\). Each Taylor coefficient is thus
determined by a finite calculation, even though \(S\) is
defined by an infinite series.

\section{Symmetry and Periodicity}

\subsection{Odd symmetry}

The exponent \(2n+1\) is odd, and the sine function is odd.
Thus
\[
a_n(-x)=-a_n(x)
\]
for every \(n\) and every real \(x\).

Summing the series gives
\[
S(-x)=-S(x).
\]

Hence \(S\) is an odd function. In particular, every even
derivative of \(S\) at the origin vanishes:
\[
S^{(2k)}(0)=0,
\qquad k=0,1,2,\ldots.
\]

This agrees with the Taylor expansion in
Equation~\eqref{eq:taylor}, which contains only odd powers.

\subsection{Periodicity of the partial sums}

Although the full series is defined on the entire real line,
each finite partial sum is periodic.

\begin{proposition}
For every nonnegative integer \(N\), the partial sum \(S_N\)
has period \(2\pi(N+1)!\).
\end{proposition}

\begin{proof}
For \(0\leq n\leq N\), replacing \(x\) by
\(x+2\pi(N+1)!\) changes the argument of the sine by
\[
\frac{2\pi(N+1)!}{(n+1)!}.
\]
Because \((n+1)!\) divides \((N+1)!\), this quantity is
an integer multiple of \(2\pi\).

The sine factor in each term is therefore unchanged.
Every term of \(S_N\) has the stated period, so their sum
has that period as well.
\end{proof}

The periods supplied by this proposition increase with
\(N\). The proposition establishes periodicity of every
finite approximation, but it does not establish that the
infinite sum has a finite period.

\section{Numerical Approximation}

\subsection{A uniform error estimate}

The partial sums provide natural approximations to \(S(x)\).
Their accuracy can be bounded independently of the value
of \(x\).

\begin{theorem}
For every nonnegative integer \(N\) and every
\(x\in\mathbb{R}\),
\[
|S(x)-S_N(x)|
\leq
\sum_{n=N+1}^{\infty}
\frac{1}{e^n(n!)^2}.
\]
Moreover,
\[
|S(x)-S_N(x)|
\leq
\frac{1}{e^{N+1}((N+1)!)^2}
\frac{1}{1-\dfrac{1}{e(N+2)^2}}.
\]
\end{theorem}

\begin{proof}
The first inequality follows directly from the triangle
inequality and the bound
\[
|a_n(x)|\leq\frac{1}{e^n(n!)^2}.
\]

Let
\[
b_n=\frac{1}{e^n(n!)^2}.
\]
For \(n\geq N+1\),
\[
\frac{b_{n+1}}{b_n}
=\frac{1}{e(n+1)^2}
\leq\frac{1}{e(N+2)^2}.
\]

Set
\[
r_N=\frac{1}{e(N+2)^2}.
\]
Since \(0<r_N<1\), comparison with a geometric series gives
\[
\sum_{n=N+1}^{\infty}b_n
\leq b_{N+1}\sum_{k=0}^{\infty}r_N^k
=\frac{b_{N+1}}{1-r_N}.
\]
Substitution of \(b_{N+1}\) proves the second inequality.
\end{proof}

This estimate gives a rigorous stopping criterion for
numerical calculations. Once the upper bound is smaller
than a chosen tolerance, the partial sum approximates
\(S(x)\) to within that tolerance for every real \(x\).

The bound is conservative because it replaces every sine
factor by its maximum possible absolute value. For a
particular input, the actual error may be substantially
smaller.

\section{Possible Generalizations}

The structure of the series suggests several variations.
For example, one may introduce a parameter \(\lambda>0\)
and consider
\[
S_{\lambda}(x)=
\sum_{n=0}^{\infty}
\frac{(-1)^n}{\lambda^n(n!)^2}
\left[
\sin\left(\frac{x}{(n+1)!}\right)
\right]^{2n+1}.
\]

For each fixed \(\lambda>0\), the same comparison argument
proves absolute and uniform convergence on \(\mathbb{R}\).
Indeed,
\[
|a_{n,\lambda}(x)|
\leq\frac{1}{\lambda^n(n!)^2},
\]
and the ratio of consecutive terms in the majorant is
\[
\frac{1}{\lambda(n+1)^2}\longrightarrow 0.
\]

Thus, the basic convergence result does not depend on
the particular choice \(\lambda=e\). Other positive
parameters can be studied using the same method.

A second variation replaces the factorial in the sine
argument by a more general sequence:
\[
T(x)=
\sum_{n=0}^{\infty}
c_n\left[\sin\left(\frac{x}{q_n}\right)\right]^{m_n},
\]
where \(c_n\), \(q_n\), and \(m_n\) are chosen sequences.
If
\[
\sum_{n=0}^{\infty}|c_n|<\infty,
\]
then \(|\sin t|\leq 1\) immediately gives absolute and
uniform convergence on the real line, provided the
remaining terms have the stated form.

This condition is sufficient rather than necessary.
If the coefficients are not absolutely summable, the
decay of the sine factors may still ensure convergence,
but a different argument would be needed.

These generalizations show that the convergence proof
belongs to a broader class of series. A more detailed
study could examine which choices of coefficients,
arguments, and exponents lead to stronger properties
or more useful approximations.

\section{Conclusion}

Scottlevoy's Factorial Series defines a bounded, odd,
real-analytic function on the real line. The Weierstrass
\(M\)-test establishes absolute and uniform convergence,
and an extension to complex inputs shows that the function
is entire. In particular, it has continuous derivatives
of every order.

The Taylor expansion describes the behavior near the
origin, where \(S(0)=0\) and \(S'(0)=1\). A global bound
controls the size of the function, while an explicit
remainder estimate controls the error of its finite
partial sums. The partial sums are also periodic, although
their periods increase with the truncation index.

Several questions remain open within the scope of this
study. These include determining whether the infinite sum
has a finite period, finding sharper bounds for its
derivatives, and investigating how the function changes
under different coefficient sequences. These questions
provide ways to extend the analysis beyond the basic
convergence and regularity results established here.


\section{Acknowledgments}
The author is a young independent student researcher and thanks the teachers and mentors who encouraged the pursuit of original mathematical exploration.

\begin{thebibliography}{9}

\bibitem{conway}
J. B. Conway,
\emph{Functions of One Complex Variable I},
2nd ed., Graduate Texts in Mathematics, vol. 11,
Springer-Verlag, New York, 1978.

\bibitem{rudin}
W. Rudin,
\emph{Principles of Mathematical Analysis},
3rd ed., International Series in Pure and Applied Mathematics,
McGraw-Hill, New York, 1976.

\end{thebibliography}

\end{document}